\subsubsection{Product rules (generalized distributive laws)}

The equality (\ref{eq.fps.prod.binary.prod-inf}) is an instance of a
\emph{product rule} -- a statement of the form \textquotedblleft a product of
sums can be expanded into one big sum\textquotedblright. The simplest product
rules are the distributive laws $a\left(  b+c\right)  =ab+ac$ and $\left(
a+b\right)  c=ac+bc$ (here, one of the sums being multiplied is a one-addend
sum); one of the next-simplest is $\left(  a+b\right)  \left(  c+d\right)
=ac+ad+bc+bd$. As far as finite sums and finite products are concerned, the
following product rule is one of the most general:\footnote{Keep in mind that
an empty Cartesian product (i.e., a Cartesian product of $0$ sets) is always a
$1$-element set; its only element is the $0$-tuple $\left(  {}\right)  $.
Thus, a sum ranging over an empty Cartesian product has exactly $1$ addend.}

\begin{proposition}
\label{prop.fps.prodrule-fin-fin}Let $L$ be a commutative ring. For every
$n\in\mathbb{N}$, let $\left[  n\right]  $ denote the set $\left\{
1,2,\ldots,n\right\}  $.

Let $n\in\mathbb{N}$. For every $i\in\left[  n\right]  $, let $p_{i,1}%
,p_{i,2},\ldots,p_{i,m_{i}}$ be finitely many elements of $L$. Then,%
\begin{equation}
\prod_{i=1}^{n}\ \ \sum_{k=1}^{m_{i}}p_{i,k}=\sum_{\left(  k_{1},k_{2}%
,\ldots,k_{n}\right)  \in\left[  m_{1}\right]  \times\left[  m_{2}\right]
\times\cdots\times\left[  m_{n}\right]  }\ \ \prod_{i=1}^{n}p_{i,k_{i}}.
\label{eq.lem.fps.prodrule-fin-fin.eq}%
\end{equation}

\end{proposition}

We can rewrite (\ref{eq.lem.fps.prodrule-fin-fin.eq}) in a less abstract way
as follows:%
\begin{align*}
&  \left(  p_{1,1}+p_{1,2}+\cdots+p_{1,m_{1}}\right)  \left(  p_{2,1}%
+p_{2,2}+\cdots+p_{2,m_{2}}\right)  \cdots\left(  p_{n,1}+p_{n,2}%
+\cdots+p_{n,m_{n}}\right) \\
&  =p_{1,1}p_{2,1}\cdots p_{n,1}+p_{1,1}p_{2,1}\cdots p_{n-1,1}p_{n,2}%
+\cdots+p_{1,m_{1}}p_{2,m_{2}}\cdots p_{n,m_{n}},
\end{align*}
where the right hand side is the sum of all $m_{1}m_{2}\cdots m_{n}$ many ways
to multiply one addend from each of the factors on the left hand side.

See \cite[solution to Exercise 6.9]{detnotes} for a formal proof of
Proposition \ref{prop.fps.prodrule-fin-fin}. (The idea is to reduce it to the
case $n=2$ by induction, then to use the discrete Fubini rule.)

Let us now move on to product rules for infinite sums and products. First, let
us extend Proposition \ref{prop.fps.prodrule-fin-fin} to a finite product of
infinite sums (which are now required to be in $K\left[  \left[  x\right]
\right]  $ in order to have a notion of summability):

\begin{proposition}
\label{prop.fps.prodrule-fin-inf}For every $n\in\mathbb{N}$, let $\left[
n\right]  $ denote the set $\left\{  1,2,\ldots,n\right\}  $.

Let $n\in\mathbb{N}$. For every $i\in\left[  n\right]  $, let $\left(
p_{i,k}\right)  _{k\in S_{i}}$ be a summable family of elements of $K\left[
\left[  x\right]  \right]  $. Then,%
\begin{equation}
\prod_{i=1}^{n}\ \ \sum_{k\in S_{i}}p_{i,k}=\sum_{\left(  k_{1},k_{2}%
,\ldots,k_{n}\right)  \in S_{1}\times S_{2}\times\cdots\times S_{n}}%
\ \ \prod_{i=1}^{n}p_{i,k_{i}}. \label{eq.lem.fps.prodrule-fin-inf.eq}%
\end{equation}
In particular, the family $\left(  \prod_{i=1}^{n}p_{i,k_{i}}\right)
_{\left(  k_{1},k_{2},\ldots,k_{n}\right)  \in S_{1}\times S_{2}\times
\cdots\times S_{n}}$ is summable.
\end{proposition}

\begin{proof}
Same method as for Proposition \ref{prop.fps.prodrule-fin-fin}, but now using
the discrete Fubini rule for infinite sums.
\end{proof}

Proposition \ref{prop.fps.prodrule-fin-inf} is rather general, but is only
concerned with \textbf{finite} products. Thus, it cannot be directly used to
justify (\ref{eq.fps.prod.binary.prod-inf}), since the product in
(\ref{eq.fps.prod.binary.prod-inf}) is infinite. Thus, we need a product rule
for infinite products of sums. Such rules are subtle and require particular
care: Not only do our sums have to be summable and our product multipliable,
but we also must avoid cases like $\left(  1-1\right)  \left(  1-1\right)
\left(  1-1\right)  \cdots$, which would produce non-summable infinite sums
when expanded (despite being multipliable). We also need on come clear about
what addends we get when we expand our products: For example, when expanding
the product
\[
\left(  1+a_{0}\right)  \left(  1+a_{1}\right)  \left(  1+a_{2}\right)
\left(  1+a_{3}\right)  \left(  1+a_{4}\right)  \cdots,
\]
we should get addends like $a_{0}\cdot1\cdot a_{2}\cdot\underbrace{1\cdot
1\cdot1\cdot\cdots}_{\text{infinitely many }1\text{s}}$ or $1\cdot a_{1}\cdot
a_{2}\cdot1\cdot a_{4}\cdot\underbrace{1\cdot1\cdot1\cdot\cdots}%
_{\text{infinitely many }1\text{s}}$, but not addends like $a_{0}\cdot
a_{1}\cdot a_{2}\cdot a_{3}\cdot\cdots$; otherwise, the right hand side of
(\ref{eq.fps.prod.binary.prod-inf}) would have to include infinite products of
$a_{i}$'s. To filter out the latter kind of addends, let us define the notion
of \textquotedblleft essentially finite\textquotedblright\ sequences or families:

\begin{definition}
\label{def.fps.prodrule.ess-fin}\textbf{(a)} A sequence $\left(  k_{1}%
,k_{2},k_{3},\ldots\right)  $ is said to be \emph{essentially finite} if all
but finitely many $i\in\left\{  1,2,3,\ldots\right\}  $ satisfy $k_{i}=0$.
\medskip

\textbf{(b)} A family $\left(  k_{i}\right)  _{i\in I}$ is said to be
\emph{essentially finite} if all but finitely many $i\in I$ satisfy $k_{i}=0$.
\end{definition}

For example, the sequence $\left(  2,4,1,0,0,0,0,\ldots\right)  $ is
essentially finite, whereas the sequence $\left(  0,1,0,1,0,1,\ldots\right)  $
(which alternates between $0$s and $1$s) is not.

Of course, Definition \ref{def.fps.prodrule.ess-fin} \textbf{(a)} is a
particular case of Definition \ref{def.fps.prodrule.ess-fin} \textbf{(b)},
since a sequence $\left(  k_{1},k_{2},k_{3},\ldots\right)  $ is the same as a
family $\left(  k_{i}\right)  _{i\in\left\{  1,2,3,\ldots\right\}  }$ indexed
by the positive integers. We also remark that Definition
\ref{def.infsum.essfin} \textbf{(a)} is a particular case of Definition
\ref{def.fps.prodrule.ess-fin} \textbf{(b)}, since a family of elements of $K$
is one particular type of family.

Now, we can finally state a version of the product rule for infinite products
of potentially infinite sums. This will help us derive
(\ref{eq.fps.prod.binary.prod-inf}) (even though the sums being multiplied in
(\ref{eq.fps.prod.binary.prod-inf}) are finite).

\begin{proposition}
\label{prop.fps.prodrule-inf-infN}Let $S_{1},S_{2},S_{3},\ldots$ be infinitely
many sets that all contain the number $0$. Set%
\[
\overline{S}=\left\{  \left(  i,k\right)  \ \mid\ i\in\left\{  1,2,3,\ldots
\right\}  \text{ and }k\in S_{i}\text{ and }k\neq0\right\}  .
\]


For any $i\in\left\{  1,2,3,\ldots\right\}  $ and any $k\in S_{i}$, let
$p_{i,k}$ be an element of $K\left[  \left[  x\right]  \right]  $. Assume
that
\begin{equation}
p_{i,0}=1\ \ \ \ \ \ \ \ \ \ \text{for any }i\in\left\{  1,2,3,\ldots\right\}
. \label{eq.prop.fps.prodrule-inf-infN.pi0=1}%
\end{equation}
Assume further that the family $\left(  p_{i,k}\right)  _{\left(  i,k\right)
\in\overline{S}}$ is summable. Then, the product $\prod_{i=1}^{\infty}%
\ \ \sum_{k\in S_{i}}p_{i,k}$ is well-defined (i.e., the family $\left(
p_{i,k}\right)  _{k\in S_{i}}$ is summable for each $i\in\left\{
1,2,3,\ldots\right\}  $, and the family $\left(  \sum_{k\in S_{i}}%
p_{i,k}\right)  _{i\in\left\{  1,2,3,\ldots\right\}  }$ is multipliable), and
we have%
\begin{equation}
\prod_{i=1}^{\infty}\ \ \sum_{k\in S_{i}}p_{i,k}=\sum_{\substack{\left(
k_{1},k_{2},k_{3},\ldots\right)  \in S_{1}\times S_{2}\times S_{3}\times
\cdots\\\text{is essentially finite}}}\ \ \prod_{i=1}^{\infty}p_{i,k_{i}}.
\label{eq.prop.fps.prodrule-inf-infN.eq}%
\end{equation}
In particular, the family $\left(  \prod_{i=1}^{\infty}p_{i,k_{i}}\right)
_{\left(  k_{1},k_{2},k_{3},\ldots\right)  \in S_{1}\times S_{2}\times
S_{3}\times\cdots\text{ is essentially finite}}$ is summable.
\end{proposition}

Note that the assumption (\ref{eq.prop.fps.prodrule-inf-infN.pi0=1}) in
Proposition \ref{prop.fps.prodrule-inf-infN} ensures that each of the sums
being multiplied contains an addend that equals $1$ (and that this addend is
named $p_{i,0}$; but this clearly does not restrict the generality of the
proposition). The equality (\ref{eq.prop.fps.prodrule-inf-infN.eq}) shows how
we can expand an infinite product of such sums. The result is a huge sum of
products (the right hand side of (\ref{eq.prop.fps.prodrule-inf-infN.eq}));
each of these products is formed by picking an addend from each sum
$\sum_{k\in S_{i}}p_{i,k}$ (just as in the finite case). The picking has to be
done in such a way that the addend $1$ gets picked from all but finitely many
of our sums (for instance, when expanding $\left(  1+x^{1}\right)  \left(
1+x^{2}\right)  \left(  1+x^{3}\right)  \cdots$, we don't want to pick the
$x^{i}$'s from all sums, as this would lead to $x^{1}x^{2}x^{3}\cdots
=$\textquotedblleft$x^{\infty}$\textquotedblright); this is why the sum on the
right hand side of (\ref{eq.prop.fps.prodrule-inf-infN.eq}) is ranging only
over the \textbf{essentially finite} sequences $\left(  k_{1},k_{2}%
,k_{3},\ldots\right)  \in S_{1}\times S_{2}\times S_{3}\times\cdots$. Finally,
the condition that the family $\left(  p_{i,k}\right)  _{\left(  i,k\right)
\in\overline{S}}$ be summable is our way of ruling out cases like $\left(
1-1\right)  \left(  1-1\right)  \left(  1-1\right)  \cdots$, which cannot be
expanded. Proposition \ref{prop.fps.prodrule-inf-infN} is not the most general
version of the product rule, but it is sufficient for many of our needs (and
we will see some more general versions below).

The rest of this subsection should probably be skipped at first reading -- it
is largely a succession of technical arguments about finite and infinite sets
in service of making rigorous what is already intuitively clear.

Before we sketch a proof of Proposition \ref{prop.fps.prodrule-inf-infN}, let
us show how (\ref{eq.fps.prod.binary.prod-inf}) can be derived from it:

\begin{proof}
[Proof of (\ref{eq.fps.prod.binary.prod-inf}) using Proposition
\ref{prop.fps.prodrule-inf-infN}.]Let $\left(  a_{0},a_{1},a_{2}%
,\ldots\right)  =\left(  a_{n}\right)  _{n\in\mathbb{N}}$ be a summable
sequence of FPSs in $K\left[  \left[  x\right]  \right]  $. We must prove the
equality (\ref{eq.fps.prod.binary.prod-inf}).

Set $S_{i}=\left\{  0,1\right\}  $ for each $i\in\left\{  1,2,3,\ldots
\right\}  $. Thus,%
\[
S_{1}\times S_{2}\times S_{3}\times\cdots=\left\{  0,1\right\}  \times\left\{
0,1\right\}  \times\left\{  0,1\right\}  \times\cdots=\left\{  0,1\right\}
^{\infty}.
\]


Define the set $\overline{S}$ as in Proposition
\ref{prop.fps.prodrule-inf-infN}. Then, $\overline{S}=\left\{  \left(
1,1\right)  ,\ \left(  2,1\right)  ,\ \left(  3,1\right)  ,\ \ldots\right\}
$. Now, set $p_{i,k}=a_{i-1}^{k}$ for each $i\in\left\{  1,2,3,\ldots\right\}
$ and each $k\in\left\{  0,1\right\}  $. Thus, the family $\left(
p_{i,k}\right)  _{\left(  i,k\right)  \in\overline{S}}$ is
summable\footnote{Indeed, this family can be rewritten as $\left(
p_{1,1},p_{2,1},p_{3,1},\ldots\right)  =\left(  a_{0}^{1},a_{1}^{1},a_{2}%
^{1},\ldots\right)  =\left(  a_{0},a_{1},a_{2},\ldots\right)  $, but we have
assumed that the latter family is summable.}, and the statement
(\ref{eq.prop.fps.prodrule-inf-infN.pi0=1}) holds as well (indeed, we have
$p_{i,0}=a_{i-1}^{0}=1$ for any $i\in\left\{  1,2,3,\ldots\right\}  $). Hence,
Proposition \ref{prop.fps.prodrule-inf-infN} can be applied.

Moreover, for each $i\in\left\{  1,2,3,\ldots\right\}  $, we have
$S_{i}=\left\{  0,1\right\}  $. Thus, for each $i\in\left\{  1,2,3,\ldots
\right\}  $, we have%
\[
\sum_{k\in S_{i}}p_{i,k}=\sum_{k\in\left\{  0,1\right\}  }p_{i,k}%
=\underbrace{p_{i,0}}_{=a_{i-1}^{0}=1}+\underbrace{p_{i,1}}_{=a_{i-1}%
^{1}=a_{i-1}}=1+a_{i-1}.
\]
Hence,%
\[
\prod_{i=1}^{\infty}\ \ \underbrace{\sum_{k\in S_{i}}p_{i,k}}_{=1+a_{i-1}%
}=\prod_{i=1}^{\infty}\left(  1+a_{i-1}\right)  =\prod_{i\in\mathbb{N}}\left(
1+a_{i}\right)
\]
(here, we have substituted $i$ for $i-1$ in the product). Thus,%
\begin{align}
\prod_{i\in\mathbb{N}}\left(  1+a_{i}\right)   &  =\prod_{i=1}^{\infty
}\ \ \sum_{k\in S_{i}}p_{i,k}\nonumber\\
&  =\underbrace{\sum_{\substack{\left(  k_{1},k_{2},k_{3},\ldots\right)  \in
S_{1}\times S_{2}\times S_{3}\times\cdots\\\text{is essentially finite}}%
}}_{\substack{=\sum_{\substack{\left(  k_{1},k_{2},k_{3},\ldots\right)
\in\left\{  0,1\right\}  ^{\infty}\\\text{is essentially finite}%
}}\\\text{(since }S_{1}\times S_{2}\times S_{3}\times\cdots=\left\{
0,1\right\}  ^{\infty}\text{)}}}\ \ \prod_{i=1}^{\infty}\underbrace{p_{i,k_{i}%
}}_{=a_{i-1}^{k_{i}}}\ \ \ \ \ \ \ \ \ \ \left(  \text{by
(\ref{eq.prop.fps.prodrule-inf-infN.eq})}\right) \nonumber\\
&  =\sum_{\substack{\left(  k_{1},k_{2},k_{3},\ldots\right)  \in\left\{
0,1\right\}  ^{\infty}\\\text{is essentially finite}}}\ \ \prod_{i=1}^{\infty
}a_{i-1}^{k_{i}}=\sum_{\substack{\left(  k_{0},k_{1},k_{2},\ldots\right)
\in\left\{  0,1\right\}  ^{\infty}\\\text{is essentially finite}%
}}\ \ \underbrace{\prod_{i=1}^{\infty}a_{i-1}^{k_{i-1}}}_{=\prod
_{i\in\mathbb{N}}a_{i}^{k_{i}}}\nonumber\\
&  \ \ \ \ \ \ \ \ \ \ \ \ \ \ \ \ \ \ \ \ \left(
\begin{array}
[c]{c}%
\text{here, we have renamed the summation}\\
\text{index }\left(  k_{1},k_{2},k_{3},\ldots\right)  \text{ as }\left(
k_{0},k_{1},k_{2},\ldots\right)
\end{array}
\right) \nonumber\\
&  =\sum_{\substack{\left(  k_{0},k_{1},k_{2},\ldots\right)  \in\left\{
0,1\right\}  ^{\infty}\\\text{is essentially finite}}}\ \ \prod_{i\in
\mathbb{N}}a_{i}^{k_{i}}. \label{pf.eq.fps.prod.binary.prod-inf.5}%
\end{align}


Now, an essentially finite sequence $\left(  k_{0},k_{1},k_{2},\ldots\right)
\in\left\{  0,1\right\}  ^{\infty}$ is just an infinite sequence of $0$s and
$1$s that contains only finitely many $1$s. Thus, such a sequence is uniquely
determined by the positions of the $1$s in it, and these positions form a
finite set, so they can be uniquely labeled as $i_{1},i_{2},\ldots,i_{k}$ with
$i_{1}<i_{2}<\cdots<i_{k}$. To put this more formally: There is a bijection%
\begin{align*}
&  \text{from }\left\{  \text{essentially finite sequences }\left(
k_{0},k_{1},k_{2},\ldots\right)  \in\left\{  0,1\right\}  ^{\infty}\right\} \\
&  \text{to }\left\{  \text{finite lists }\left(  i_{1},i_{2},\ldots
,i_{k}\right)  \text{ of nonnegative integers such that }i_{1}<i_{2}%
<\cdots<i_{k}\right\}
\end{align*}
that sends each sequence $\left(  k_{0},k_{1},k_{2},\ldots\right)  \in\left\{
0,1\right\}  ^{\infty}$ to the list of all $i\in\mathbb{N}$ satisfying
$k_{i}=1$ (written in increasing order). Furthermore, if this bijection sends
an essentially finite sequence $\left(  k_{0},k_{1},k_{2},\ldots\right)
\in\left\{  0,1\right\}  ^{\infty}$ to a finite list $\left(  i_{1}%
,i_{2},\ldots,i_{k}\right)  $, then
\begin{align*}
\prod_{i\in\mathbb{N}}a_{i}^{k_{i}}  &  =\left(  \prod_{\substack{i\in
\mathbb{N};\\k_{i}=0}}\ \ \underbrace{a_{i}^{k_{i}}}_{=a_{i}^{0}=1}\right)
\left(  \prod_{\substack{i\in\mathbb{N};\\k_{i}=1}}\ \ \underbrace{a_{i}%
^{k_{i}}}_{=a_{i}^{1}=a_{i}}\right)  \ \ \ \ \ \ \ \ \ \ \left(  \text{since
each }k_{i}\text{ is either }0\text{ or }1\right) \\
&  =\underbrace{\left(  \prod_{\substack{i\in\mathbb{N};\\k_{i}=0}}1\right)
}_{=1}\left(  \prod_{\substack{i\in\mathbb{N};\\k_{i}=1}}a_{i}\right)
=\prod_{\substack{i\in\mathbb{N};\\k_{i}=1}}a_{i}=a_{i_{1}}a_{i_{2}}\cdots
a_{i_{k}}%
\end{align*}
(since $\left(  i_{1},i_{2},\ldots,i_{k}\right)  $ is the list of all
$i\in\mathbb{N}$ satisfying $k_{i}=1$). Thus, we can use this bijection to
re-index the sum on the right hand side of
(\ref{pf.eq.fps.prod.binary.prod-inf.5}); we obtain%
\[
\sum_{\substack{\left(  k_{0},k_{1},k_{2},\ldots\right)  \in\left\{
0,1\right\}  ^{\infty}\\\text{is essentially finite}}}\ \ \prod_{i\in
\mathbb{N}}a_{i}^{k_{i}}=\sum_{i_{1}<i_{2}<\cdots<i_{k}}a_{i_{1}}a_{i_{2}%
}\cdots a_{i_{k}}.
\]
Thus, (\ref{pf.eq.fps.prod.binary.prod-inf.5}) rewrites as%
\[
\prod_{i\in\mathbb{N}}\left(  1+a_{i}\right)  =\sum_{i_{1}<i_{2}<\cdots<i_{k}%
}a_{i_{1}}a_{i_{2}}\cdots a_{i_{k}}.
\]
This proves (\ref{eq.fps.prod.binary.prod-inf}).
\end{proof}

We note that (\ref{eq.fps.prod.binary.prod-inf}) can be rewritten as%
\begin{equation}
\prod_{i\in\mathbb{N}}\left(  1+a_{i}\right)  =\sum_{\substack{J\text{ is a
finite}\\\text{subset of }\mathbb{N}}}\ \ \prod_{i\in J}a_{i}.
\label{eq.fps.prod.binary.prod-inf2}%
\end{equation}
Indeed, the right hand side of this equality is precisely the right hand side
of (\ref{eq.fps.prod.binary.prod-inf}).

We will prove Proposition \ref{prop.fps.prodrule-inf-infN} soon. First,
however, let us generalize Proposition \ref{prop.fps.prodrule-inf-infN} to
products indexed by arbitrary sets instead of $\left\{  1,2,3,\ldots\right\}
$:

\begin{proposition}
\label{prop.fps.prodrule-inf-inf}Let $I$ be a set. For any $i\in I$, let
$S_{i}$ be a set that contains the number $0$. Set%
\[
\overline{S}=\left\{  \left(  i,k\right)  \ \mid\ i\in I\text{ and }k\in
S_{i}\text{ and }k\neq0\right\}  .
\]


For any $i\in I$ and any $k\in S_{i}$, let $p_{i,k}$ be an element of
$K\left[  \left[  x\right]  \right]  $. Assume that
\begin{equation}
p_{i,0}=1\ \ \ \ \ \ \ \ \ \ \text{for any }i\in I.
\label{eq.prop.fps.prodrule-inf-inf.pi0=1}%
\end{equation}
Assume further that the family $\left(  p_{i,k}\right)  _{\left(  i,k\right)
\in\overline{S}}$ is summable. Then, the product $\prod_{i\in I}\ \ \sum_{k\in
S_{i}}p_{i,k}$ is well-defined (i.e., the family $\left(  p_{i,k}\right)
_{k\in S_{i}}$ is summable for each $i\in I$, and the family $\left(
\sum_{k\in S_{i}}p_{i,k}\right)  _{i\in I}$ is multipliable), and we have%
\begin{equation}
\prod_{i\in I}\ \ \sum_{k\in S_{i}}p_{i,k}=\sum_{\substack{\left(
k_{i}\right)  _{i\in I}\in\prod_{i\in I}S_{i}\\\text{is essentially finite}%
}}\ \ \prod_{i\in I}p_{i,k_{i}}. \label{eq.prop.fps.prodrule-inf-inf.eq}%
\end{equation}
In particular, the family $\left(  \prod_{i\in I}p_{i,k_{i}}\right)  _{\left(
k_{i}\right)  _{i\in I}\in\prod_{i\in I}S_{i}\text{ is essentially finite}}$
is summable.
\end{proposition}

Clearly, Proposition \ref{prop.fps.prodrule-inf-infN} is the particular case
of Proposition \ref{prop.fps.prodrule-inf-inf} for $I=\left\{  1,2,3,\ldots
\right\}  $. The proof of Proposition \ref{prop.fps.prodrule-inf-inf} is a
long-winded reduction to the finite case; nothing substantial is going on in
it. Thus, I recommend skipping it unless specifically interested. For the sake
of convenience, before proving Proposition \ref{prop.fps.prodrule-inf-inf},
let me restate Proposition \ref{prop.fps.prodrule-fin-inf} in a form that
makes it particularly easy to use:

\begin{proposition}
\label{prop.fps.prodrule-fin-infJ}Let $N$ be a finite set. For any $i\in N$,
let $\left(  p_{i,k}\right)  _{k\in S_{i}}$ be a summable family of elements
of $K\left[  \left[  x\right]  \right]  $. Then,%
\begin{equation}
\prod_{i\in N}\ \ \sum_{k\in S_{i}}p_{i,k}=\sum_{\left(  k_{i}\right)  _{i\in
N}\in\prod_{i\in N}S_{i}}\ \ \prod_{i\in N}p_{i,k_{i}}.
\label{eq.lem.fps.prodrule-fin-infJ.eq}%
\end{equation}
In particular, the family $\left(  \prod_{i\in N}p_{i,k_{i}}\right)  _{\left(
k_{i}\right)  _{i\in N}\in\prod_{i\in N}S_{i}}$ is summable.
\end{proposition}

\begin{proof}
[Proof of Proposition \ref{prop.fps.prodrule-fin-infJ}.]This is just
Proposition \ref{prop.fps.prodrule-fin-inf}, with the indexing set $\left\{
1,2,\ldots,n\right\}  $ replaced by $N$. It can be proved by reindexing the
products (or directly by induction on $\left\vert N\right\vert $).
\end{proof}

Let me furthermore state an infinite version of Lemma
\ref{lem.fps.prod.irlv.fin}, which will be used in the proof of Proposition
\ref{prop.fps.prodrule-inf-inf} given below:

\begin{lemma}
\label{lem.fps.prod.irlv.inf}Let $a\in K\left[  \left[  x\right]  \right]  $
be an FPS. Let $\left(  f_{i}\right)  _{i\in J}\in K\left[  \left[  x\right]
\right]  ^{J}$ be a summable family of FPSs. Let $n\in\mathbb{N}$. Assume that
each $i\in J$ satisfies
\begin{equation}
\left[  x^{m}\right]  \left(  f_{i}\right)  =0\ \ \ \ \ \ \ \ \ \ \text{for
each }m\in\left\{  0,1,\ldots,n\right\}  .
\label{eq.lem.fps.prod.irlv.inf.ass}%
\end{equation}
Then,
\[
\left[  x^{m}\right]  \left(  a\prod_{i\in J}\left(  1+f_{i}\right)  \right)
=\left[  x^{m}\right]  a\ \ \ \ \ \ \ \ \ \ \text{for each }m\in\left\{
0,1,\ldots,n\right\}  .
\]

\end{lemma}

\begin{proof}
[Proof of Lemma \ref{lem.fps.prod.irlv.inf} (sketched).]The idea here is to
argue that the first $n+1$ coefficients of $a\prod_{i\in J}\left(
1+f_{i}\right)  $ agree with those of $a\prod_{i\in M}\left(  1+f_{i}\right)
$ for some finite subset $M$ of $J$. Then, apply Lemma
\ref{lem.fps.prod.irlv.fin} to this subset $M$. The details of this argument
can be found in Section \ref{sec.details.gf.prod}.
\end{proof}

With this lemma proved, we have all necessary prerequisites for the proof of
Proposition \ref{prop.fps.prodrule-inf-inf}. The proof, however, is rather
long due to the bookkeeping required, and therefore has been banished to the
appendix (Section \ref{sec.details.gf.prod}, to be specific).

We notice that Proposition \ref{prop.fps.prodrule-inf-inf} (and thus also
Proposition \ref{prop.fps.prodrule-inf-infN}) can be generalized slightly by
lifting the requirement that all sets $S_{i}$ contain $0$ (this means that the
sums being multiplied no longer need to contain $1$ as an addend). See
Exercise \ref{exe.fps.prodrule-inf-inf-0} for this generalization. This lets
us expand products such as $x\left(  1+x^{1}\right)  \left(  1+x^{2}\right)
\left(  1+x^{3}\right)  \cdots$ (but of course, we could just as well expand
this particular product by splitting off the $x$ factor and applying
Proposition \ref{prop.fps.prodrule-inf-infN}).

\subsubsection{Another example}

As another bit of practice with infinite products of FPSs, let us prove the
following identity by Euler (\cite[\S 326]{Euler48}):

\begin{proposition}
[Euler]\label{prop.gf.prod.euler-odd}We have%
\[
\prod_{i>0}\left(  1-x^{2i-1}\right)  ^{-1}=\prod_{k>0}\left(  1+x^{k}\right)
.
\]

\end{proposition}

\begin{proof}
First, it is easy to see that the families $\left(  1+x^{k}\right)  _{k>0}$
and $\left(  1-x^{k}\right)  _{k>0}$ and $\left(  1-x^{2k}\right)  _{k>0}$ and
$\left(  1-x^{2i-1}\right)  _{i>0}$ are multipliable\footnote{Indeed, this
follows from Theorem \ref{thm.fps.1+f-mulable}, since the families $\left(
x^{k}\right)  _{k>0}$ and $\left(  -x^{k}\right)  _{k>0}$ and $\left(
-x^{2k}\right)  _{k>0}$ and $\left(  -x^{2i-1}\right)  _{i>0}$ are summable.}.
This shows that the product on the right hand side of Proposition
\ref{prop.gf.prod.euler-odd} is well-defined.

Proposition \ref{prop.fps.div-mulable} \textbf{(a)} shows that the family
$\left(  \dfrac{1}{1-x^{2i-1}}\right)  _{i>0}$ is multipliable (since the
families $\left(  1\right)  _{i>0}$ and $\left(  1-x^{2i-1}\right)  _{i>0}$
are multipliable, and since the FPS $1-x^{2i-1}$ is invertible for each
$i>0$). In other words, the family $\left(  \left(  1-x^{2i-1}\right)
^{-1}\right)  _{i>0}$ is multipliable. Thus, the product on the left hand side
of Proposition \ref{prop.gf.prod.euler-odd} is well-defined.

For each $k>0$, we have $1+x^{k}=\dfrac{1-x^{2k}}{1-x^{k}}$ (since
$1-x^{2k}=1-\left(  x^{k}\right)  ^{2}=\left(  1-x^{k}\right)  \left(
1+x^{k}\right)  $). Thus,%
\begin{align*}
\prod_{k>0}\underbrace{\left(  1+x^{k}\right)  }_{=\dfrac{1-x^{2k}}{1-x^{k}}}
&  =\prod_{k>0}\dfrac{1-x^{2k}}{1-x^{k}}\\
&  =\dfrac{\prod_{k>0}\left(  1-x^{2k}\right)  }{\prod_{k>0}\left(
1-x^{k}\right)  }\ \ \ \ \ \ \ \ \ \ \left(  \text{by Proposition
\ref{prop.fps.div-mulable} \textbf{(b)}}\right) \\
&  =\dfrac{\left(  1-x^{2}\right)  \left(  1-x^{4}\right)  \left(
1-x^{6}\right)  \left(  1-x^{8}\right)  \cdots}{\left(  1-x^{1}\right)
\left(  1-x^{2}\right)  \left(  1-x^{3}\right)  \left(  1-x^{4}\right)
\cdots}\\
&  =\dfrac{1}{\left(  1-x^{1}\right)  \left(  1-x^{3}\right)  \left(
1-x^{5}\right)  \left(  1-x^{7}\right)  \cdots}\\
&  \ \ \ \ \ \ \ \ \ \ \ \ \ \ \ \ \ \ \ \ \left(
\begin{array}
[c]{c}%
\text{since all }1-x^{i}\text{ factors with }i\text{ even}\\
\text{cancel out (see below for details)}%
\end{array}
\right) \\
&  =\dfrac{1}{1-x^{1}}\cdot\dfrac{1}{1-x^{3}}\cdot\dfrac{1}{1-x^{5}}%
\cdot\dfrac{1}{1-x^{7}}\cdot\cdots\\
&  \ \ \ \ \ \ \ \ \ \ \ \ \ \ \ \ \ \ \ \ \left(  \text{by Proposition
\ref{prop.fps.div-mulable} \textbf{(b)}}\right) \\
&  =\prod_{i>0}\dfrac{1}{1-x^{2i-1}}=\prod_{i>0}\left(  1-x^{2i-1}\right)
^{-1},
\end{align*}
which is precisely the claim of Proposition \ref{prop.gf.prod.euler-odd}.

The step where we cancelled out the $1-x^{i}$ factors with $i$ even can be
justified as follows: We know that the family $\left(  1-x^{k}\right)  _{k>0}$
as well as its subfamilies $\left(  1-x^{k}\right)  _{k>0\text{ is even}}$ and
$\left(  1-x^{k}\right)  _{k>0\text{ are odd}}$ are multipliable. Thus,
Theorem \ref{prop.fps.union-mulable} \textbf{(b)} yields%
\[
\prod_{k>0}\left(  1-x^{k}\right)  =\left(  \prod_{\substack{k>0\\\text{is
even}}}\left(  1-x^{k}\right)  \right)  \cdot\left(  \prod
_{\substack{k>0\\\text{is odd}}}\left(  1-x^{k}\right)  \right)  .
\]
In other words,%
\begin{align*}
&  \left(  1-x^{1}\right)  \left(  1-x^{2}\right)  \left(  1-x^{3}\right)
\left(  1-x^{4}\right)  \cdots\\
&  =\left(  \left(  1-x^{2}\right)  \left(  1-x^{4}\right)  \left(
1-x^{6}\right)  \left(  1-x^{8}\right)  \cdots\right) \\
&  \ \ \ \ \ \ \ \ \ \ \cdot\left(  \left(  1-x^{1}\right)  \left(
1-x^{3}\right)  \left(  1-x^{5}\right)  \left(  1-x^{7}\right)  \cdots\right)
.
\end{align*}
Hence,%
\[
\dfrac{\left(  1-x^{1}\right)  \left(  1-x^{2}\right)  \left(  1-x^{3}\right)
\left(  1-x^{4}\right)  \cdots}{\left(  1-x^{2}\right)  \left(  1-x^{4}%
\right)  \left(  1-x^{6}\right)  \left(  1-x^{8}\right)  \cdots}=\left(
1-x^{1}\right)  \left(  1-x^{3}\right)  \left(  1-x^{5}\right)  \left(
1-x^{7}\right)  \cdots.
\]
Taking reciprocals, we thus obtain%
\[
\dfrac{\left(  1-x^{2}\right)  \left(  1-x^{4}\right)  \left(  1-x^{6}\right)
\left(  1-x^{8}\right)  \cdots}{\left(  1-x^{1}\right)  \left(  1-x^{2}%
\right)  \left(  1-x^{3}\right)  \left(  1-x^{4}\right)  \cdots}=\dfrac
{1}{\left(  1-x^{1}\right)  \left(  1-x^{3}\right)  \left(  1-x^{5}\right)
\left(  1-x^{7}\right)  \cdots},
\]
just as we claimed. Thus, the proof of Proposition
\ref{prop.gf.prod.euler-odd} is complete.
\end{proof}

Let us try to interpret Proposition \ref{prop.gf.prod.euler-odd}
combinatorially, by expanding both sides.

First, we use (\ref{eq.fps.prod.binary.prod-inf}) to expand the right hand
side:%
\begin{align}
\prod_{k>0}\left(  1+x^{k}\right)   &  =\left(  1+x^{1}\right)  \left(
1+x^{2}\right)  \left(  1+x^{3}\right)  \left(  1+x^{4}\right)  \cdots
\nonumber\\
&  =\sum_{\substack{i_{1},i_{2},\ldots,i_{k}\in\left\{  1,2,3,\ldots\right\}
;\\i_{1}<i_{2}<\cdots<i_{k}}}\underbrace{x^{i_{1}}x^{i_{2}}\cdots x^{i_{k}}%
}_{=x^{i_{1}+i_{2}+\cdots+i_{k}}}\ \ \ \ \ \ \ \ \ \ \left(  \text{by
(\ref{eq.fps.prod.binary.prod-inf})}\right) \nonumber\\
&  =\sum_{\substack{i_{1},i_{2},\ldots,i_{k}\in\left\{  1,2,3,\ldots\right\}
;\\i_{1}<i_{2}<\cdots<i_{k}}}x^{i_{1}+i_{2}+\cdots+i_{k}}\nonumber\\
&  =\sum_{n\in\mathbb{N}}d_{n}x^{n}, \label{eq.prop.gf.prod.euler-odd.d}%
\end{align}
where $d_{n}$ is the \# of all strictly increasing tuples $\left(  i_{1}%
<i_{2}<\cdots<i_{k}\right)  $ of positive integers such that $n=i_{1}%
+i_{2}+\cdots+i_{k}$. We can rewrite this definition as follows: $d_{n}$ is
the \# of ways to write $n$ as a sum of distinct positive integers, with no
regard for the order\footnote{\textquotedblleft No regard for the
order\textquotedblright\ means that, for example, $3+4+1$ and $1+3+4$ count as
the same way of writing $8$ as a sum of distinct integers.}.

Next, let us expand the left hand side: For each positive integer $i$, we have%
\begin{align*}
\left(  1-x^{2i-1}\right)  ^{-1}  &  =1+x^{2i-1}+\left(  x^{2i-1}\right)
^{2}+\left(  x^{2i-1}\right)  ^{3}+\cdots\\
&  \ \ \ \ \ \ \ \ \ \ \ \ \ \ \ \ \ \ \ \ \left(
\begin{array}
[c]{c}%
\text{by substituting }x^{2i-1}\text{ for }x\text{ in the}\\
\text{equality }\left(  1-x\right)  ^{-1}=1+x+x^{2}+x^{3}+\cdots
\end{array}
\right) \\
&  =1+x^{2i-1}+x^{2\left(  2i-1\right)  }+x^{3\left(  2i-1\right)  }%
+\cdots=\sum_{k\in\mathbb{N}}x^{k\left(  2i-1\right)  }.
\end{align*}
Multiplying these equalities over all $i>0$, we find%
\begin{align}
&  \prod_{i>0}\left(  1-x^{2i-1}\right)  ^{-1}\nonumber\\
&  =\prod_{i>0}\ \ \sum_{k\in\mathbb{N}}x^{k\left(  2i-1\right)  }\nonumber\\
&  =\sum_{\substack{\left(  k_{1},k_{2},k_{3},\ldots\right)  \in
\mathbb{N}^{\infty}\\\text{is essentially finite}}}\underbrace{\prod
_{i>0}x^{k_{i}\left(  2i-1\right)  }}_{\substack{=x^{k_{1}\cdot1}x^{k_{2}%
\cdot3}x^{k_{3}\cdot5}\cdots\\=x^{k_{1}\cdot1+k_{2}\cdot3+k_{3}\cdot5+\cdots}%
}}\ \ \ \ \ \ \ \ \ \ \left(
\begin{array}
[c]{c}%
\text{by Proposition \ref{prop.fps.prodrule-inf-infN},}\\
\text{applied to }S_{i}=\mathbb{N}\text{ and }p_{i,k}=x^{k\left(  2i-1\right)
}%
\end{array}
\right) \nonumber\\
&  =\sum_{\substack{\left(  k_{1},k_{2},k_{3},\ldots\right)  \in
\mathbb{N}^{\infty}\\\text{is essentially finite}}}x^{k_{1}\cdot1+k_{2}%
\cdot3+k_{3}\cdot5+\cdots}\nonumber\\
&  =\sum_{n\in\mathbb{N}}o_{n}x^{n}, \label{eq.prop.gf.prod.euler-odd.o}%
\end{align}
where $o_{n}$ is the \# of all essentially finite sequences $\left(
k_{1},k_{2},k_{3},\ldots\right)  \in\mathbb{N}^{\infty}$ such that $k_{1}%
\cdot1+k_{2}\cdot3+k_{3}\cdot5+\cdots=n$. I claim that $o_{n}$ is the \# of
ways to write $n$ as a sum of (not necessarily distinct) odd positive
integers, without regard to the order. (Why? Because if we can write $n$ as a
sum of $k_{1}$ many $1$s, $k_{2}$ many $3$s, $k_{3}$ many $5$s and so on, then
$\left(  k_{1},k_{2},k_{3},\ldots\right)  \in\mathbb{N}^{\infty}$ is an
essentially finite sequence such that $k_{1}\cdot1+k_{2}\cdot3+k_{3}%
\cdot5+\cdots=n$; and conversely, if $\left(  k_{1},k_{2},k_{3},\ldots\right)
\in\mathbb{N}^{\infty}$ is such a sequence, then $n$ can be written as a sum
of $k_{1}$ many $1$s, $k_{2}$ many $3$s, $k_{3}$ many $5$s and so on.)

Proposition \ref{prop.gf.prod.euler-odd} tells us that $\prod_{k>0}\left(
1+x^{k}\right)  =\prod_{i>0}\left(  1-x^{2i-1}\right)  ^{-1}$. In view of
(\ref{eq.prop.gf.prod.euler-odd.d}) and (\ref{eq.prop.gf.prod.euler-odd.o}),
we can rewrite this as
\[
\sum_{n\in\mathbb{N}}d_{n}x^{n}=\sum_{n\in\mathbb{N}}o_{n}x^{n}.
\]
Therefore, $d_{n}=o_{n}$ for each $n\in\mathbb{N}$. Thus, we have proven the
following purely combinatorial statement:

\begin{theorem}
[Euler]\label{thm.gf.prod.euler-comb}Let $n\in\mathbb{N}$. Then, $d_{n}=o_{n}%
$, where

\begin{itemize}
\item $d_{n}$ is the \# of ways to write $n$ as a sum of distinct positive
integers, without regard to the order;

\item $o_{n}$ is the \# of ways to write $n$ as a sum of (not necessarily
distinct) odd positive integers, without regard to the order.
\end{itemize}
\end{theorem}

\begin{example}
Let $n=6$. Then, $d_{n}=4$, because the ways to write $n=6$ as a sum of
distinct positive integers are%
\[
6,\ \ \ \ \ \ \ \ \ \ 1+5,\ \ \ \ \ \ \ \ \ \ 2+4,\ \ \ \ \ \ \ \ \ \ 1+2+3
\]
(don't forget the first of these ways, trivial as it looks!). On the other
hand, $o_{n}=4$, because the ways to write $n=6$ as a sum of odd positive
integers are%
\[
1+5,\ \ \ \ \ \ \ \ \ \ 3+3,\ \ \ \ \ \ \ \ \ \ 3+1+1+1,\ \ \ \ \ \ \ \ \ \ 1+1+1+1+1+1.
\]

\end{example}

We will soon learn a bijective proof of Theorem \ref{thm.gf.prod.euler-comb}
as well (see the Second proof of Theorem \ref{thm.pars.odd-dist-equal} below).

\subsubsection{Infinite products and substitution}

Proposition \ref{prop.fps.subs.rules} \textbf{(h)} has an analogue for
products instead of sums:

\begin{proposition}
\label{prop.fps.subs.rule-infprod}If $\left(  f_{i}\right)  _{i\in I}\in
K\left[  \left[  x\right]  \right]  ^{I}$ is a multipliable family of FPSs,
and if $g\in K\left[  \left[  x\right]  \right]  $ is an FPS satisfying
$\left[  x^{0}\right]  g=0$, then the family $\left(  f_{i}\circ g\right)
_{i\in I}\in K\left[  \left[  x\right]  \right]  ^{I}$ is multipliable as well
and we have $\left(  \prod_{i\in I}f_{i}\right)  \circ g=\prod_{i\in I}\left(
f_{i}\circ g\right)  $.
\end{proposition}

\begin{proof}
[Proof of Proposition \ref{prop.fps.subs.rule-infprod} (sketched).]See Section
\ref{sec.details.gf.prod}.
\end{proof}

\subsubsection{Exponentials, logarithms and infinite products}

Recall Definition \ref{def.fps.Exp-Log-maps}. As we saw in Lemma
\ref{lem.fps.Exp-Log-additive}, the exponential map $\operatorname*{Exp}$ and
the logarithm map $\operatorname*{Log}$ can be used to convert sums into
products and vice versa (under certain conditions). This also extends to
infinite sums and infinite products:

\begin{proposition}
\label{prop.fps.Exp-Log-infsum}Assume that $K$ is a commutative $\mathbb{Q}%
$-algebra. Let $\left(  f_{i}\right)  _{i\in I}\in K\left[  \left[  x\right]
\right]  ^{I}$ be a summable family of FPSs in $K\left[  \left[  x\right]
\right]  _{0}$. Then, $\left(  \operatorname*{Exp}f_{i}\right)  _{i\in I}$ is
a multipliable family of FPS in $K\left[  \left[  x\right]  \right]  _{1}$,
and we have $\sum_{i\in I}f_{i}\in K\left[  \left[  x\right]  \right]  _{0}$
and%
\[
\operatorname*{Exp}\left(  \sum_{i\in I}f_{i}\right)  =\prod_{i\in I}\left(
\operatorname*{Exp}f_{i}\right)  .
\]

\end{proposition}

\begin{proposition}
\label{prop.fps.Exp-Log-infprod}Assume that $K$ is a commutative $\mathbb{Q}%
$-algebra. Let $\left(  f_{i}\right)  _{i\in I}\in K\left[  \left[  x\right]
\right]  ^{I}$ be a multipliable family of FPSs in $K\left[  \left[  x\right]
\right]  _{1}$. Then, $\left(  \operatorname*{Log}f_{i}\right)  _{i\in I}$ is
a summable family of FPS in $K\left[  \left[  x\right]  \right]  _{0}$, and we
have $\prod\limits_{i\in I}f_{i}\in K\left[  \left[  x\right]  \right]  _{1}$
and%
\[
\operatorname*{Log}\left(  \prod\limits_{i\in I}f_{i}\right)  =\sum_{i\in
I}\left(  \operatorname*{Log}f_{i}\right)  .
\]

\end{proposition}

We leave the proofs of these two propositions as an exercise (Exercise
\ref{exe.fps.Exp-Log-infsum}).
